\begin{textsample}{NEG Dim 1 – al – Score 4.00 – al/al\_ef\_22.txt}  \label{ex:f1_neg_134}
O Eixo da Oralidade diz respeito às práticas de linguagem que ocorrem em situação oral, seja face a face ou por outros modos de interação, tais como : aula dialogada, webconferência, mensagem gravada, spot de campanha, jingle, seminário, debate, programa de rádio, entrevista, declamação de poemas ( com ou sem efeitos sonoros ), peça teatral, apresentação de cantigas e canções, playlist comentada de músicas, vlog de game, contação de histórias, diferentes tipos de podcasts e vídeos, dentre outras.
O tratamento das práticas orais propostas na BNCC e contempladas também neste referencial compreende : - Consideração e reflexão sobre as condições de produção dos textos orais que regem a circulação de diferentes gêneros nas diferentes mídias e campos de atividade humana - Refletir sobre diferentes contextos e situações sociais em que se produzem textos orais e sobre as diferenças em termos formais, estilísticos e linguísticos que esses contextos determinam, incluindo -se aí a multimodalidade e a multissemiose. - Conhecer e refletir sobre as tradições orais e seus gêneros, considerando -se as práticas sociais em que tais textos surgem e se perpetuam, bem como os sentidos que geram. - Compreensão de textos orais - Proceder a uma escuta ativa, voltada para questões relativas ao contexto de produção dos textos, para o conteúdo em questão, para a observação de estratégias discursivas e dos recursos linguísticos e multissemióticos mobilizados, bem como dos elementos paralinguísticos e cinésicos. - Produção de textos orais - Produzir textos pertencentes a gêneros orais diversos, considerando -se aspectos relativos ao planejamento, à produção, ao redesign, à avaliação das práticas realizadas em situações de interação social específicas. - Compreensão dos efeitos de sentidos provocados pelos usos de recursos linguísticos e multissemióticos em textos pertencentes a gêneros diversos - Identificar e analisar efeitos de sentido decorrentes de escolhas de volume, timbre, intensidade, pausas, ritmo, efeitos sonoros, sincronização, expressividade, gestualidade etc.
Produzir textos levando em conta efeitos possíveis. - Relação entre fala e escrita - Estabelecer relação entre fala e escrita, levando -se em conta o modo como as duas modalidades se articulam em diferentes gêneros e práticas de linguagem ( como jornal de TV, programa de rádio, apresentação de seminário, mensagem instantânea etc. ), as semelhanças e as diferenças entre modos de falar e de registrar a escrita e os aspectos sociodiscursivos, composicionais e linguísticos de cada modalidade sempre relacionados com os gêneros em questão. - Oralizar o texto escrito, considerando -se as situações sociais em que tal tipo de atividade acontece, seus elementos paralinguísticos e cinésicos, dentre outros. - Refletir sobre as variedades linguísticas, adequando sua produção a esse contexto.

% matched lemmas: 
\end{textsample}
